\odiseachapter{Séptima parte}{Cíclope}

Hasta aquí, Christopher y el que suscribe no han llegado a las manos. Darle un tajo a Helena en el rostro es simplista, pero la escena de Esparta fluye, Lupita Nyong'o está divina y Jon Bernthal interpreta a un convincente Menelao. Eliminar a Nausícaa es un pecado capital y quitar de en medio a los feacios una blasfemia digna de excomunión, pero, ¡ay!, así son las exigencias del guion.

En cambio, eliminar la más célebre ---y característica--- de todas las astucias de Odiseo le va a costar a Nolan penar en el séptimo círculo del infierno, junto con Michael Bay (por \emph{Transformers}), Roland Emmerich (\emph{Independence Day}) y Santiago Segura, por habernos hecho reír a pesar de todos los pesares de \emph{Torrente}.

Pero dejemos que sea Homero quien nos guíe en este fenomenal episodio. Recordemos que Ulises está contando sus cuitas en la corte de Alcínoo ---por tanto narrando en primera persona---. Han salido por piernas después de vérselas con los cícones y de pasar por el país de los lotófagos ---un pequeño episodio que Nolan se salta, a cambio de convertir a Charlize, digo a Calipso, en una traficante del género---. Lo cierto es que en mi versión de Ítaca, el viejo Ulises recuerda, quizás inexactamente ---o quizás es el Poeta el que se salta detalles--- lo que pasó:

\begin{verse}
Tal vez nunca creíste en dioses,\\
y por esta razón siempre estabas cauteloso y alerta.\\
¿Recuerdas el país de los comedores de loto? Todos tus hombres\\
estaban drogados y roncando cinco minutos después de desembarcar,\\
tú no. De alguna manera, sentías el peligro\\
en las bellas personas que te ofrecían flores,\\
peligro en el cielo azul, en el agua tranquila,\\
peligro en el día perfecto, que invitaba al viajero\\
a tumbarse en la arena, a cerrar sus ojos cansados\\
y dormir.\\
Tú, dormir no lo hiciste. La belleza\\
te asustaba más que el único y ardiente ojo de Polifemo.\\
Esos cuerpos perfectos y ágiles, desnudos bajo el sol,\\
de alguna manera deletreaban Muerte.

Belleza y muerte. Cuando caminas por la costa\\
de tu pequeña isla, escuchando los gritos de las gaviotas,\\
---a menudo oyes las voces de Áyax, de Aquiles,\\
incluso de Héctor, llamándote desde el Hades---,\\
cuando paseas entre los naranjos, oliendo la fragancia\\
de las flores de primavera ---tan dulces como el loto que rehusaste comer---,\\
cuando miras el cielo estrellado durante la noche,\\
¿por qué tu alma se estremece de anhelo y de deseo?\\
¿por qué tu corazón se siente pesado y ligero a la vez?\\
Tal vez sea la noción\\
de que gaviotas y flores y estrellas\\
aún seguirán ahí,\\
cuando tú te hayas ido.
\end{verse}

Por fin llegan a la isla de los Cíclopes, que resultan ser una panda de gandules, dándose la gran vida a expensas de los favores de los dioses:

\begin{verse}
De allí otra vez navegamos con el corazón sombritriste,\\
hasta que a patria de unos altivos sin ley, de los Cíclopes,\\
fuimos a dar, que fiando ellos siempre en los dioses felices,\\
ni aran la tierra ni siembran semilla que luego germine,\\
pero sin siembra ni arada a ellos todo les nace y les rinde,\\
los trigos y las cebadas y los racimos en vides\\
llenos de vino que Zeus con su lluvia les infla y bendice.\\
Y leyes ellos no tienen, ni plaza en la que se decide,\\
que es en las crestas de altas montañas en donde viven\\
dentro de cuéncavas grutas; cada uno sus leyes prescribe\\
para sus hijos y esposas, sin que unos de otros se cuiden.
\end{verse}

Unos «altivos sin ley» que fían en los dioses felices para que les llenen la despensa y la bota de vino, nada menos. ¿Y por qué? ¿Qué razones tienen los dioses para favorecerlos? Homero no lo explica, pero más tarde el anciano Ulises aventurará una respuesta. Por el momento, sin embargo, el aventurero que acaba de pegarle fuego a la ciudad de los cícones se dispone a ver qué puede saquear aquí.

\begin{verse}
A un lado y fuera del puerto hay una isla pequeña, alargada,\\
que del país de los Cíclopes cerca no está ni apartada,\\
y es toda bosque; por ella salvajes se mueven las cabras\\
innumerables, pues no las ahuyentan veredas humanas\\
ni entran allí cazadores que ronden a la emboscada\\
y se castiguen siguiéndoles rastro, alto allá en la montaña.
\end{verse}

Ulises cuenta que sus naves embarrancan, con fortuna, en las playas de la isla:

\begin{verse}
Allí, pues, fueron las naves, y a fe que algún dios nos guiaba\\
esa negrísima noche, que a ver nada allí se nos daba.\\
Solo una niebla cerrada en redor, y ni luna asomaba\\
para nosotros del cielo, las nubes no la dejaban.\\
Nadie la isla allí vio con sus ojos, ni vimos las largas\\
olas tampoco rodar a lo largo por sobre la playa\\
hasta que en tierra encallamos las naves bellibancadas.\\
Ellas varadas, echamos las velas abajo a las tablas\\
y al fin nos desembarcamos sobre el arenal de las aguas.\\
Nos entregamos al sueño allí hasta las claras del alba.
\end{verse}

Y al despertarse, los aqueos se encuentran el festín servido, como quien dice:

\begin{verse}
Cuando sus dedos de rosa la Aurora asomó tempranera,\\
fuimos a andar por la isla, mirándolo todo en sorpresa.\\
Las ninfas, hijas de Zeus cabrillante, al pasar nos ojean\\
cabras montesas, y a bien, por aviarle a mi gente la cena.\\
Y ya los arcos recorvos y las azagayas punteras\\
de nuestros barcos sacamos, y en tres dividimos las fuerzas\\
y disparamos; y a gana y contento nos dio un dios las presas.\\
Doce las naves conmigo; tocaron cada una de ellas\\
a nueve cabras, y diez me apartaron a mí en preferencia.\\
Fue así un estar todo el día, hasta hundirse ya el sol, como en fiesta\\
saboreando con muy dulce vino tasajos a espuertas;\\
que no se había agotado el tinto de nuestras bodegas,\\
que aún nos quedaba del que trasvasamos en ánforas llenas\\
cuando la santa ciudad de los cícones plaza fue nuestra.\\
Viendo allí estábamos tierra de Cíclopes ---era bien cerca---\\
y humo, y sintiendo sus voces, y las de sus cabras y ovejas.\\
Y cuando el sol se hundió y vino a caer la tiniebla,\\
allí a dormir nos tendimos al borde del mar por la arena.
\end{verse}

No pueden quejarse: se han dado un festín y han repuesto fuerzas. La prudencia aconsejaría seguir camino. Pero Ulises no sería Ulises si no fuera curioso (y rapaz) y al día siguiente decide explorar, a ver si puede llevarse algo más:

\begin{verse}
Cuando sus dedos de rosa la Aurora asomó tempranera,\\
los reuní y para todos hablé porque todos me oyeran:\\
«Ahora quedaos, mis fieles, aquí sin moveros del sitio;\\
que yo me voy a acercar con mi nave y mi marinerío\\
para saber de esos hombres y ver quiénes son, de cuál signo,\\
si son soberbios, brutales y de la justicia no amigos,\\
o acogedores y bien temerosos del cielo en su espíritu».
\end{verse}

Ya le vale, al saqueador de ciudades, ponerse tan estupendo con los habitantes de la isla, que a duras penas van a resultar más «soberbios, brutales y de la justicia no amigos» que él y su tripulación. Pero si los dioses pueden permitirse una doble moral, tampoco se la niega Homero a nuestro héroe milmañas:

\begin{verse}
Esto diciendo, subí yo a la nave, ordenando a mi tropa\\
que se embarcara también y soltara los cables de popa.\\
Y a esa mi voz embarcaron y junto a los remos se apostan,\\
y, uno tras otro sentados, golpeaban la mar canizosa.\\
En cuanto a aquella región arribamos ---se hallaba muy próxima---,\\
vimos allá, en un extremo, una cueva mirando a las olas\\
bien en lo alto, y envuelta en laureles. De noche era choza\\
de un apretado rebaño de cabras y ovejas; de rocas\\
bien incrustadas en tierra se alzaba una cerca a redonda,\\
con largos pinos entre ellas y encinas altifrondosas.\\
De noche allí se encovaba un hombre gigante, que, a solas,\\
apacentaba sus hatos aparte, sin ver a persona, y\\
lejos, muy lejos de todos, vivía ignorante de normas.
\end{verse}

A Homero le encantan los paralelismos. Hemos visto ya varios y veremos más. Aquí hay uno que salta a la vista. Calipso y el Cíclope viven ambos en cuevas, lo que sugiere, dicho sea de paso, que el Poeta no está dando a entender que ese tipo de morada corresponda a un habitante rudo o miserable, sino, quizás, a un ser con ciertos atributos sobrenaturales. Tampoco dice que el Cíclope sea un monstruo, ni que tenga un solo ojo\footnote{Ese detalle no aparece en todo el canto, aunque se puede deducir por el contexto que si Polifemo tiene dos, solo ve por uno de ellos. El asunto es importante en mi propio poema. Quien le da un solo ojo a los cíclopes es Hesíodo, en su \emph{Teogonía}.}. Eso sí, nos avisa de que es un gigante, y nada le da pie a pensar que vaya a recibirlos amablemente, a pesar de lo cual Odiseo sigue empeñado en la visita:

\begin{verse}
Entonces orden les di a mis fieles correligionarios\\
de que esperaran allí y tuvieran cuidado del barco,\\
y, tras coger doce hombres, con ellos marché, los más válidos,\\
y un vino negro en un odre de cabra, oloroso, regalo\\
que tuvo a bien a mí hacerme Marón, el hijo de Evanto\\
y sacerdote de Apolo ---que es dios de la guarda de Ismaro---,\\
por dar a su hijo y su esposa y a él nuestro abrazo y cuidado
\end{verse}

El vino en cuestión, nos explica, es tan fuerte que para beberlo hay que mezclar una medida del licor con veinte de agua.

\begin{verse}
Y cada vez que este rojo licor melidulce probaron,\\
veinte las partes de agua que echaban para un solo vaso,\\
y era tan dulce el olor que exhalaba el caldero al echarlo,\\
tan de los dioses, que nadie querría estar de él alejado.\\
De este llevé un gran pellejo, y avíos también en un saco,\\
que muy enseguida a mi espíritu ardido le vino el cuidado\\
de que, de fuerza muy grande, en nada saldríame al paso\\
hombre salvaje que a ley y justicia iba a serles extraño.
\end{verse}

No hay otro episodio de la \emph{Odisea} donde veamos más en directo la astucia de Ulises. Se lleva a cuestas un gran pellejo de ese orujo capaz de emborrachar a un ejército porque se barrunta que el gigante de la cueva no tiene nada de amable y educado. ¿Por qué se empeña en ir a visitarlo entonces? La respuesta, conociendo a Odiseo, es más que obvia. Quiere ver qué puede sacarle, por las buenas ---la famosa \emph{xenía}--- o, aunque no lo dice, por las malas, cosa de la que ya sabemos que es muy capaz.

El relato continúa:

\begin{verse}
Rápidamente alcanzamos la cueva, pero a él allí dentro\\
no lo encontramos, que andaba en el pasto un hatajo apaciendo.\\
Así que por la espelunca mirábamos esto y aquello:\\
zarzos cargados de quesos, apriscos de chivos repletos\\
y de borregos, y bien separados por tiempo entre ellos:\\
ahí los primeros, allá los medianos, y aparte los nuevos.\\
Todas las cántaras de lecherías nadaban en suero,\\
cuencos y buenas colodras en las que juntaba el ordeño.\\
Y aquí ya entonces hablaron mis hombres, y su único ruego\\
era marcharnos cuanto antes después de coger unos quesos,\\
y regresar para hacer de borregos y chivos arreo\\
desde el redil a la nave, y al fin poner agua por medio.\\
Ah, pero yo no hice caso ---¡y cuánto mejor haberlo hecho!---\\
por verlo a él y, quizás, por ganar, como huésped, su obsequio.\\
Su aparición bienvenida no fue para mis compañeros.
\end{verse}

¡Et voilà! Para los estándares de la época, Polifemo vive en la cueva del tesoro, que Homero nos describe con todo detalle. Quesos, chivos, borregos, cántaros de leche. Los hombres que ha traído consigo, sensatamente, proponen hacerse con unos quesos y salir por piernas antes de que llegue el gigante. Pero Ulises, haciendo gala de valentía ---pero también de una imprudencia infantil, vista la talla del cíclope, y de una arrogancia que le impide pensar en su tripulación--- insiste en quedarse, mitad por curiosidad, mitad por avaricia («Ah, pero yo no hice caso ---¡y cuánto mejor haberlo hecho!--- por verlo a él y, quizás, por ganar, como huésped, su obsequio»).

Así que se instalan sin que nadie les invite, se comen un queso que no es suyo y esperan, como tontos, a que aparezca el dueño de la cueva. En cuanto este asoma, como era de esperar, se llevan un susto de muerte:

\begin{verse}
Así que hicimos un fuego, quemamos en él las primicias,\\
cogimos queso, comimos, y dentro esperamos sin prisa\\
a que llegara el pastor. Un brazado imponente traía\\
de leña seca para prepararse después la comida.\\
Él con estrépito al suelo la echó, dentro ya de la sima,\\
y, estremecidos, al fondo corrimos de sus galerías.
\end{verse}

Polifemo mete las ovejas que piensa ordeñar en la cueva y sella la puerta con una piedra gigante.

\begin{verse}
De esa su gorda ovejada en la gran espelunca metía\\
solo la que iba a ordeñar, los machos afuera tenía,\\
cabrones junto a carneros, guardándose en la corraliza.\\
Luego empotró, levantándola a pulso, una piedra grandísima,\\
ponderosísima; ni veintidós carretas macizas\\
de cuatro ruedas, y a una, del suelo la apalancarían;\\
tanto y tan recio el peñasco que puerta a la entrada le hacía.
\end{verse}

Como buen pastor, se sienta a ordeñar su rebaño de manera aplicada y sistemática, lo que no da pie a considerarlo un bárbaro:

\begin{verse}
Y se sentó ya a ordeñar, todo tal como manda doctrina,\\
cabras y ovejas, y púsole luego a cada una su cría.\\
Cuajó después la mitad de la blanca leche y, deprisa,\\
cama le dio amontonándola en las mimbreurdidas cestillas;\\
la otra mitad a unos cuencos la echó por tener su bebida\\
a mano siempre, y con ella la cena la acompañaría.
\end{verse}

Cuando ha terminado la faena, aviva el fuego y repara en Ulises y su gente. Les dice:

\begin{verse}
«Extraños, ¿quiénes ahí sois? ¿Desde dónde la aguosa derrota?\\
¿Vais por negocios por caso o a lo que delante se os ponga,\\
como piratas que de un lado a otro barriendo las olas\\
vida se juegan llevándole el mal a la tierra que tocan?».
\end{verse}

Otro ejemplo de que Homero no da puntada sin hilo. En el canto III, cuando Telémaco visita a Néstor en Pilos, este los recibe amablemente ---sin saber aún quiénes son--- como exigen las leyes de la hospitalidad y después de obsequiarles con comida y bebida les hace la misma pregunta:

\begin{verse}
Y cuando avientan por fin de bebida y comida el deseo,\\
Néstor gerenio es primero en hablarles, buencaballero:\\
«Ya es más gentil, ahora sí, con los huéspedes ir a otras cosas\\
y preguntar quiénes son, pues contentos ya están de bucólica.\\
Huéspedes, pues, ¿quiénes sois? ¿Desde dónde esta aguosa derrota?\\
¿Vais por negocios por caso o a lo que delante se os ponga,\\
como piratas que de un lado a otro barriendo las olas\\
vida se juegan llevándole el mal a la tierra que tocan?».
\end{verse}

A lo que Telémaco contesta identificándose como el hijo de Odiseo. Se trata de una fórmula ritual entre el anfitrión que ofrece la \emph{xenía} y el huésped que se beneficia de esta. Un detalle sutil es que Tobal traduce \gr{ὦ ξεῖνοι} por «huéspedes» en casa de Néstor y por «extraños» en la cueva del cíclope. Se da el caso de que ambas traducciones son correctas: en griego clásico, «huésped» y «extraño» son la misma palabra. Pero Tobal puede permitirse remarcar el carácter de «invitados» en casa de Néstor, en contraposición con el de «extraños» en la cueva del cíclope, de connotación mucho más hostil. Más importante, Néstor pregunta quiénes son a sus invitados después de darles de comer, es decir, después de haberles demostrado hospitalidad. El cíclope lo hace inmediatamente, lo que no augura nada bueno. Los aqueos, con razón, se mueren de miedo, pero Ulises no pierde la calma:

\begin{verse}
Eso nos dijo, y nuestros corazones se nos resquebraban,\\
amedrentados por esa voz honda y sus trazas gigantas.\\
Tuve, aun así, yo palabras para responder su demanda:

Somos aqueos a los que de Troya hasta aquí han traído lejos\\
sobre la gran hondonal de la mar, unos y otros, los vientos;\\
vamos a casa, pero otros caminos, otros senderos\\
aquí nos tienen; será que idearlo ha querido así Zeus.\\
De Agamenón el Atrida tenemos a bien ser ejército;\\
gloria mayor que la suya ahora bajo los cielos\\
no la hay: ciudad tanta fue la que hundió, tantos fueron sus muertos.
\end{verse}

Dicho lo cual, le suplica hospitalidad invocando el mandamiento de los dioses. Eso sí, la súplica incluye una amenaza implícita, ya que Zeus se toma muy a mal que un potencial anfitrión se salte la \emph{xenía}:

\begin{verse}
A tus rodillas venimos ahora, en este tu suelo,\\
por si los dones de huésped llegáramos a merecerlos,\\
u otro presente nos dieras, de huéspedes mandamiento.\\
Respeta tú, buen amigo, a los dioses; venimos en ruego, y\\
del suplicante y del huésped es Zeus vengador justiciero,\\
Zeus Hospedal, que acompaña a los huéspedes bienreverendos.
\end{verse}

Pero a Polifemo las leyes de Zeus le traen sin cuidado:

\begin{verse}
Dije así, y él respondió desde su corazón despiadado:\\
«O eres un bobo, extranjero, o de un pago llegas lejano,\\
cuando me pides que tema a los dioses y evite enfadarlos.\\
Que miramiento en los Cíclopes no hay ni por Zeus cabrilámpigo\\
ni por sus dioses felices, pues somos más fuertes, y aun harto.\\
No guardaría de mal, por el odio de Zeus evitarlo,\\
ni a ti ni a tus compañeros si no lo mandase mi ánimo.\\
En fin… ¿Y dónde, al venir, has dejado tu bien hecho barco?,\\
te lo pregunto por curiosidad, ¿está ahí o en el cabo?».
\end{verse}

Veamos. El amigo Christopher se pasa toda su película obsesionado con el desastre que supone violar la \emph{xenía}, pero si examinamos con cierta atención sus argumentos, veremos que están sujetos con pinzas. Muy cierto, el caballo de Troya fue una treta astuta y no poco vil, pero después de todo los troyanos y los aqueos estaban en guerra y nadie les mandaba a los de Ilión meter el equinomaderoso en su alcázar y acostarse a dormir tan campantes. No hay en esa situación un extranjero que pide hospitalidad y un anfitrión que se la niega. Y de hecho, si leemos atentamente la \emph{Odisea}, veremos que afirma,  a lo largo de sus muchos cantos, exactamente lo contrario de lo que dice Nolan. Néstor y Menelao respetan la \emph{xenía}, como también lo hace Calipso y, si me apuras, Circe, aunque a la hechicera trenzadivina hay que echarle de comer aparte. Y por supuesto, los feacios se definen como la encarnación misma de la hospitalidad. Los dos ejemplos en los que la \emph{xenía} se viola son el caso de los pretendientes ---donde los huéspedes abusan del anfitrión--- y el caso de Polifemo, que no solo se niega a ser anfitrión y se ríe de los dioses, sino que agrede a sus huéspedes. Y nuestro admirado director, con la misma poca fortuna con que prescinde del ejemplo más vivo de hospitalidad en la \emph{Odisea}, ignora el mejor ejemplo en el que esta se viola de manera racional y consciente. El problema del cíclope de Nolan es que es un pobre hombre, un simple, que se come a los aqueos como quien se come una sardina en escabeche, sin darle, aparentemente, importancia al hecho, mientras que el Polifemo de Odiseo es artero, cruel, arrogante... Y astuto. Fijémonos en su última pregunta. «Por curiosidad, ¿dónde has dejado tus naves?». No solo piensa comerse a los rehenes que tiene encerrados en su cueva, sino que pretende enterarse de dónde está el resto de la tripulación y con seguridad que no es para hacerles una visita de cortesía.

Y aquí se repite otro patrón que ya conocemos, el del enfrentamiento entre dos enemigos inteligentes y arteros. Ulises ya se las había visto con las mañas de Helena en Troya y ahora le toca responder a los ardides de Polifemo. Y si en Ilión sale victorioso usando la resistencia pasiva ---negándose a caer en la trampa de Helena y guardando silencio---, en la cueva del cíclope el mientisagaz da con la mayor de todas sus célebres tretas. Responde:

\begin{verse}
«Mi nave la destrozó Poseidón territodatremante\\
en los confines de vuestra tierra; después de acercarme\\
a un arrecife, a unas rocas me echó, con el viento empujándome.\\
Ruina escarpada esquivé yo con estos que tienes delante».
\end{verse}

Al enemigo, cuanta menos información, mejor. Polifemo se cansa de conversar y decide prepararse un aperitivo:

\begin{verse}
Dije así, y no respondió desde su corazón despiadado,\\
solo saltó, y a mis hombres echó de repente las manos,\\
y a dos cogió y, cual cachorros de perro, empezó a golpearlos\\
contra la tierra, y corrían los sesos, el suelo encharcado.\\
Miembro por miembro los despedazó, su comida aprestando.
\end{verse}

Una escena espeluznante, donde de paso aprendemos que entre los aqueos todavía no había cuajado la idea del bienestar animal. El cíclope golpea a los hombres contra el suelo «cual cachorros de perro», es decir, el trato que le dan estos a las crías no deseadas de su ---supuestamente--- mejor amigo. La descripción sigue, con toda crueldad:

\begin{verse}
Comía, y nada dejaba, igual que león somontano,\\
vísceras, carnes y huesos de tuétanos entreverados.\\
Nosotros, manos al cielo, a Zeus elevamos el llanto\\
viendo ese horror, y ya al corazón lo ganó el desamparo.\\
Cuando su estómago enorme el Cíclope lo hubo llenado\\
comiendo carne de hombre y la pura leche tragando,\\
entre sus cabras dormido quedó en la mitad de aquel antro.
\end{verse}

Ulises piensa en matarlo y se da cuenta de que si lo hace quedan atrapados en la cueva, ya que no tienen forma de mover el enorme peñasco ---aquí Nolan sigue el poema al pie de la letra---. Así que la fuerza bruta no vale y hay que recurrir a la astucia. Llega la mañana siguiente:

\begin{verse}
Cuando sus dedos de rosa la Aurora asomó tempranera,\\
hizo una lumbre y se puso a ordeñar las lustrosas ovejas\\
bien a conciencia, y su cría después a cada una le deja.\\
Luego de haber atendido con buena faena sus jeras,\\
vuelve a echar mano a otros dos de los míos y se los almuerza.\\
Desayunado, la gorda ovejada sacó de la cueva,\\
moviendo aquel gran peñasco, con facilidad, de la puerta;\\
y lo volvió a colocar, cual si tapa a la aljaba pusiera.\\
Ya monte arriba a su gorda manada entre silbos la lleva el\\
Cíclope; yo me quedé de venganza rumiando manera,\\
si se me daba vengarme y venía a mi rezo Atenea.\\
Y esta es la cuenta que me pareció de verdad la más buena.
\end{verse}

Otro detalle impresionante es el hecho de que el cruel Polifemo trata con mimo exquisito a sus ovejas. Quizás, después de todo, no es que Polifemo sea un monstruo, sino que, en su escala de valores, las vidas de los aqueos no valen nada. ¿Tiene Ulises derecho a quejarse de semejante trato después de haber arrasado Ilión? Pero claro, Ulises no se queda en la cueva ponderando la relatividad del bien y el mal, sino tramando venganzas. Y aquí viene la famosa escena del madero puntiforme.

\begin{verse}
Junto al redil dejó el Cíclope un tronco muy largo y aún tierno;\\
él lo cortó de un olivo pensando en usarlo, ya seco,\\
como cayado. Que un mástil nosotros diríamos, viéndolo,\\
era de un negro navío de carga de veinte remeros,\\
de esos que van más allá de la gran hondonal; tan soberbio\\
nos parecía ese tronco por su longitud y tan grueso.\\
Y a él me fui, y una braza de largo corté, más o menos,\\
que les llevé, porque lo desbastaran, a mis compañeros.\\
Y lo afinaron; después lo cogí y le agucé yo un extremo\\
que endurecí volteándolo bien en el vívido fuego.\\
Y lo escondí con cuidado cubriéndolo con el estiércol\\
que allí, por toda la cueva, se amontonaba en el suelo.\\
Entonces les ordené sortear entre ellos el resto\\
de los osados que levantarían conmigo el madero\\
para estregarlo en su ojo en cuanto viniérale el sueño.\\
Esos que yo elegiría, los mismos bendijo el sorteo,\\
cuatro, y yo mismo que hacía el quinto elegido con ellos.
\end{verse}

Impresionante la riqueza de detalles que da Homero. Cuenta cómo afilan el palo y luego lo esconden con precisión cinematográfica. Al atardecer, regresa el cíclope y vuelve a las andadas, pero Odiseo pone en marcha su plan:

\begin{verse}
Atardecía y volvió de apacer con su grey lanilinda.\\
De la rolliza manada en la gran espelunca metía\\
toda esta vez; no dejó animal fuera por la corraliza,\\
oliéndose algo quizás, o quizás algún dios lo movía.\\
Luego empotró, levantándola a pulso, la piedra grandísima,\\
y se sentó ya a ordeñar, todo tal como manda doctrina,\\
cabras y ovejas, poniéndole luego a cada una su cría.\\
Cuando quedaron aquellas sus jeras tan bien atendidas,\\
garra de nuevo echó a dos de mis hombres, y cena se hacía.\\
Entonces fui yo hacia el Cíclope, y dije mientras le tendía\\
desde mis manos de aquel vino negro una buena escudilla:
\end{verse}

\begin{verse}
«Cíclope, bebe, después de comer carne de hombre, este vino,\\
para que sepas de qué buen beber iba lleno el navío.\\
En libación para ti lo traía, si, compadecido,\\
me devolvías a casa; mas no hay quien aguante tus ímpetus.\\
Maldito, ¿cómo podría llegársete en lo sucesivo\\
hombre ninguno de tantos que hay, si en ti ley no hay ni tino?».
\end{verse}

\begin{verse}
Le dije; y él lo cogió y lo apuró; y disfrutó a maravilla\\
bebiendo el dulce licor, y pidió otra segunda enseguida:
\end{verse}

¡Funciona!, debió de pensar Odiseo. Polifemo acaba de beberse un vaso lleno del orujo infernal y pide más. El plan de emborracharlo va viento en popa y, de nuevo, Homero es más sutil que el bienfamoso director. En la película, el monstruo pocasluces se echa a dormir después de darse el festín y los aqueos aprovechan para cegarlo. En el original, Ulises lo emborracha primero, garantizando que lo deja fuera de juego y puede maniobrar a sus anchas. El gigante matón, excesivamente confiado, no se lo ve venir, pide más y suelta otra de esas frases donde Homero teje su historia con más finura aún que el sudario de Penélope:

\begin{verse}
«Dame otra vez, generoso, y dime ya ahora tu nombre\\
para que pueda yo darte regalo de huésped que goces;\\
también la tierra triguera a los Cíclopes deja sus dones,\\
buenos racimos para que de vino el que llueve los colme;\\
pero de néctar y de ambrosía es torrente el que escondes».
\end{verse}

Le pide el nombre a cambio de un prometido regalo. La burla no puede ser más descarada. Los feacios solo preguntarán por el nombre del héroe después de colmarlo de regalos y queda muy claro aquí que poco puede esperarse de la generosidad de este truhan comehombres. Pero Odiseo lleva un as en la manga. Acaba de emborracharlo y luego le dice cómo se llama:

\begin{verse}
Tal dijo; y ya estaba yo alargándole el vino fogueño.\\
Tres veces fui y se lo di, tres con poco saber él bebiendo.\\
Y cuando al Cíclope el vino le entró hasta los mismos adentros,\\
le dije entonces así, como adulzando mi verbo:
\end{verse}

\begin{verse}
«Cíclope, me has preguntado mi nombre glorioso; y en nada\\
te lo diré; dame tú ese regalo que has dicho que dabas.\\
Ninguno tengo por nombre; Ninguno, que así a mí me llaman\\
siempre mi madre y mi padre, y aquellos que me acompañan».
\end{verse}

Ulises dice llamarse \gr{Οὖτις} (outis), o sea Nadie, Ninguno. Polifemo, satisfecho, revela, malvadamente, que su regalo será comérselo el último.

\begin{verse}
Dije así, y él respondió desde su corazón despiadado:\\
«Me comeré yo a Ninguno el último tras sus mandados;\\
ellos irán por delante. Y eso será tu regalo».
\end{verse}

A continuación se echa a dormir y el resto ya lo sabemos. Los aqueos lo ciegan, pinchándole el ojo con el puntiagudo palo\footnote{De manera algo intrigante, Homero, que nunca ha dicho que Polifemo tenga un solo ojo, aquí da por supuesto que ese es el caso. ¿O quizás Polifemo sufre de alguna enfermedad y su segundo ojo está velado, o estrábico, o excesivamente miope?}. Una vez más, la precisión narrativa del Poeta es espeluznante:

\begin{verse}
Ellos, alzando aquel tronco de olivo, aguzado a un extremo,\\
lo hincaron bien en su ojo, y cargándome yo sobre el leño,\\
lo atornillaba como el que taladra de nave un madero\\
con la barrena, mientras por debajo tirando de cueros\\
otros la bailan de un lado y de otro en un giro perpetuo.\\
Así, apoyada en su ojo, a la estaca afilada en el fuego\\
dábamos baile; bañaba la sangre el ardiente madero.\\
Se chamuscaron sus párpados ambos y cejas del vuelo\\
de la abrasada pupila y sus fibras en chisporroteo.\\
Y como cuando la azuela o un hacha grande el herrero\\
en agua fría la sume en medio de silbos tremendos\\
por darle temple, porque ahí va a estar luego la fuerza del hierro,\\
así silbaba su ojo al roce del palo oliveño.
\end{verse}

Dudo que Flaubert fuera capaz de escribir un párrafo más preciso y más terrorífico. Polifemo se retuerce de dolor y llama, gimiendo, a los otros cíclopes ---que no aparecen en la película de Nolan---. Sus paisanos se acercan y le preguntan desde fuera de la cueva:

\begin{verse}
«¿Qué el daño tanto que a ti, Polifemo, te trae a esta queja\\
en la bendita noche inmortal, y sin sueño nos dejas?\\
¿Alguno te anda tirando, y muy maldegrado, de ovejas?\\
¿Alguno te anda matando con artería o con fuerza?».
\end{verse}

Y la respuesta de Polifemo pone de manifiesto el ardid más glorioso de la épica universal:

\begin{verse}
«Con artería y no a fuerza Ninguno, queridos, me mata».
\end{verse}

Los otros cíclopes, naturalmente, piensan que al compadre le deben doler las muelas y el asunto queda zanjado:

\begin{verse}
Y las palabras aladas que a él en respuesta le dieron:\\
«Pues si en verdad fuerza alguno no te hace y estás solo ahí dentro,\\
de enfermedad que del gran Zeus viene imposible el remedio.\\
Mira si al padre tú vas, Poseidón soberano, con rezos».
\end{verse}

Ulises, encantado de haberse conocido, nos dice:

\begin{verse}
Eso decían marchando de allí; mi alma estaba risueña\\
porque mi nombre les engañó y fue intachable la treta.
\end{verse}

El resto es cine. Polifemo retira la piedra, los aqueos se esconden entre las ovejas y se escapan de la cueva, llevándose de paso a los carneros. Han salido del apuro y todo lo que necesitan hacer ahora es salir por remos. Ah, pero Ulises no puede quedarse calladito, no puede soportar la idea de marcharse sin provocar a su enemigo:

\begin{verse}
«No ibas tú, Cíclope, no, a comerte en tu cuéncava gruta\\
los camaradas de un hombre sinfuerza con saña bien cruda.\\
Esa maldad llegaría a alcanzarte en igual desmesura,\\
maldito, que sin reparo hiciste de huéspedes gula\\
en tu morada: los dioses y Zeus se han cobrado su suma».
\end{verse}

Polifemo se enfada y les tira un peñasco:

\begin{verse}
Dije, y se encolerizó más aún, y con toda su alma,\\
tras desgajar el picacho de un alto peñón, nos lo lanza\\
justo por frente de nuestra goleta proalazulada;\\
para alcanzar el timón en su extremo faltó casi nada.
\end{verse}

A pesar de lo cual, consiguen salir del lío. ¿Escarmienta por ello el muysagaz Odiseo? ¿Le dicta la prudencia dejarse de bravuconadas? Nada de eso. Es demasiado presumido para largarse sin dejarle claro al gigante que le ha tomado el pelo.
Sus hombres, mucho más sensatos que él, le suplican que se calle:

\begin{verse}
«Terco, ¿por qué otra vez quieres airar a ese hombre salvaje,\\
él, que a un disparo en el mar nos llevó hace un momento la nave\\
de nuevo a tierra y allí ya nos vimos en último trance?\\
Basta un sonido que le hagas o que una palabra lo emplace\\
para que nuestras cabezas y palos del barco él machaque\\
lanzando un áspero meño. Que tanto tiene de alcance».
\end{verse}

El de Ítaca no les hace ni caso y bravuconea:

\begin{verse}
«Cíclope, si un día a ti se llegara otro bienmoridero y\\
te preguntara que quién te dejó ciego el ojo y tan feo,\\
que te cegó, le dirás, el cercoasolador Odiseo,\\
hijo de Laertes, que en Ítaca tiene su casa y asiento».
\end{verse}

Y, por supuesto, Polifemo pide venganza a su padre Poseidón.

\begin{verse}
«Óyeme tú, Poseidón, territodaseñor blavicrino,\\
si de verdad tuyo soy y te glorias de ser padre mío,\\
dame que no halle el cercoasolador Odiseo el camino,\\
hijo de Laertes, que en Ítaca tiene su casa y destino.\\
Mas, si pisar tierra patria y la casa techalta es el sino\\
que, a más de ver a los suyos, le viene correspondido,\\
que tarde llegue, y no bien, y ya todos sus hombres caídos,\\
y en nave de otros, y techo se encuentre de penas ahíto».
\end{verse}

Ulises, que se cree tan astuto, acaba de demostrar ser un insensato, deshaciendo su propia treta. Polifemo no podía pedirle a Poseidón ---que presumiblemente no estaba mirando en ese momento--- venganza contra Nadie, como no pudo pedirles a los cíclopes ayuda cuando Ninguno lo mataba con artería y no con fuerza. Pero el de Laertes acaba de darle los datos completos, que le hacen falta al dios del mar para perseguirlo por el resto de su viaje. ¡Y todo por orgullo, todo por arrogancia! En este encuentro entre dos enemigos despiadados, ambos demuestran ser, en última instancia, unos necios.

Pero hay otro detalle interesante. Ulises, de eso no le cabe duda a ninguna atenta lectora, a ningún perspicaz lector, es un narrador del que no nos podemos fiar. Hay algo en su historia del cíclope que no acaba de cuadrar. Su enemigo es, a la vez, demasiado listo y demasiado tonto ---¿A quién se le ocurre aceptar una y otra vez el vino que le ofrece Ulises después de devorar a varios aqueos?---; demasiado malo y demasiado bueno ---se come a los hombres de dos en dos pero mima con esmero a sus ovejas---. ¿Y si lo que nos cuenta no fuera exactamente cierto? Dado lo mentiroso que es, no sería ninguna sorpresa.

Así que quizás no nos sorprenda escuchar sus reflexiones, ya viejo, en Ítaca:

\begin{verse}
Y luego, por la noche, junto al fuego,\\
tus nietos suben a tu regazo y te piden una historia.\\
¡Polifemo! Gritan. Por favor, cuéntanos, abuelo,\\
cómo derrotaste al monstruo.

Así que repites el cuento una y otra vez,\\
en el que el héroe, todo valor e ingenio,\\
engaña a la fea bestia.

Esto es, al fin y al cabo, lo que todos quieren oír.\\
Y no solo los niños,\\
sino también amigos que vienen al palacio,\\
comerciantes de todas las islas, con sus ánforas llenas de mercancías,\\
pastores que pasan con sus rebaños,\\
sirvientes y soldados, incluso tu hijo y tu nuera,\\
a todos les gusta sentarse en torno al fuego\\
e imaginar a un Polifemo grande, feo y brutal,\\
listo para comerse a los aqueos uno detrás de otro,\\
pero demasiado estúpido para entender las estratagemas del héroe.

Y cada vez que estiras la madeja,\\
la verdad se disuelve un poco más en amables mentiras,\\
como la sal en el agua.\\
La verdad ---que todo el mundo conoció una vez,\\
y era tan obvia, en cualquier caso---: los cíclopes eran conocidos\\
por ser solo los hijos rechazados del Olimpo,\\
pobres, deformes e infelices, torturados\\
por huesos que crecían demasiado para su propio bien.\\
Sus ojos estaban enfermos ---tenían dos, como cualquiera,\\
pero a menudo uno de ellos estaba cubierto de densas cataratas---;\\
ciertamente, ese era el caso de Polifemo,\\
ese pobre y gentil muchacho algo retrasado,\\
cuyo único pecado era ser poco agraciado.
\end{verse}

El Polifemo del amigo Christopher tiene un eco de este muchacho retrasado que describe el poema\footnote{Y por si algún mal pensado se lo está preguntando, \emph{Abandonando Ítaca} (\url{https://eolasediciones.es/libro/abandonando-itaca-leaving-ithaca/}) precede en un par de años a la película de Nolan.}, lo que me viene bien para insistir en que no me molesta en absoluto que se reinterprete o incluso se reinvente el clásico ---yo mismo lo hago y hay una abundante literatura de variantes que enriquece aún más la inmortal épica---. El problema no es, en mi opinión, un cíclope estúpido, sino la falta de ambición para sacarle partido. El Polifemo de Nolan hace lo mismo que el de Homero y no aporta nada a la versión original, mucho más interesante. Pero elimina el architruco de Ulises y apunta a que este se enemista con Poseidón en un acto de justa venganza: le lanza una flecha al cíclope, provocando la ira del dios ---como si haberle sacado un ojo a su chico un rato antes no fuera suficiente--- por sus hombres asesinados, cuando el propio héroe cuenta, en la corte de Alcínoo, que lo hace por pura arrogancia.

El viejo Ulises de mi \emph{Abandonando Ítaca} sigue así:

\begin{verse}
¿Cómo es posible, te sorprende, que nadie se pregunte\\
por tus razones para detenerte en la isla?\\
Claro, tu relato argumenta que necesitabas comida y agua,\\
olvidando mencionar, sin embargo, que el cíclope ofrecía\\
provisión abundante y que aceptaba a cambio, como pago,\\
unas pocas baratijas.\\
Olvidando también, convenientemente, explicar que tú sospechabas\\
que tenía en algún lugar oro escondido. Y así, el ingenioso Ulises\\
decidió hacer una fiesta, embriagarlos a todos,\\
y robar sus cuevas mientras ellos roncaban.

Mala suerte que Polifemo no bebiera bastante,\\
mala suerte que tuviera que dar la batalla,\\
y tuvieses que usar la violencia.\\
Cegarlo no fue tan difícil como siempre pretendes,\\
erais muchos, y aunque fuerte,\\
no lo era lo bastante como para tener una oportunidad\\
contra tu banda de carniceros.\\
Ah, pero cómo lloró cuando tú\\
le clavaste esa viga en el ojo sano,\\
cómo llamó a su madre y te preguntó por qué,\\
incapaz de entender por qué su amigo Nadie le había hecho tanto daño.

Esa parte de la historia es cierta, al menos.\\
Los otros llamaron: ¿Quién te ha cegado? Y el pobre Polifemo gritó.\\
Nadie lo hizo.\\
Sus gritos te persiguen todavía,\\
después de tantos años.

Tal vez esta sea la razón por la cual, cada vez que cuentas tu historia,\\
él se vuelve más grande, y más mezquino, y le crecen dientes en el rostro\\
y un ojo solo, más rojo que la sangre que se derrama de las gargantas\\
de los troyanos que tú y Diomedes matasteis por deporte.

Y todo el mundo ríe y aplaude, y acepta felizmente el hecho de que el vencedor\\
debe de estar diciendo la verdad.
\end{verse}
